% 三4已完成.tex - 完成国家自然科学基金项目情况

申请人尚无已结题的国家自然科学基金项目。

申请人目前承担的首个国家自然科学基金项目为青年科学基金（批准号：12404568），执行期为2025年1月至2027年12月，目前处于执行第二年。以下汇报该项目的阶段性进展，说明与本申请项目的关联：

\subsection*{青年基金项目（12404568）阶段性进展}

\textbf{1. 三角晶格编码方案与自动化小工具搜索}

项目组提出了三角晶格编码方案~\cite{Pan2025Triangular}，显著提升了里德堡原子阵列的编码质量（详见立项依据与研究基础）。更重要的是，该工作开发了基于整数规划的自动化小工具搜索方法，验证了计算机辅助规约合成的可行性，为本项目LLM驱动的规模化合成奠定方法论基础。相关论文正在审稿中，并已申请里德堡原子平台问题映射的专利。

\textbf{2. 动态全息技术}

项目组发展了基于自动微分的动态全息图生成方法，已发表于\emph{Physical Review Applied}，为里德堡原子阵列实验平台提供支撑。

\textbf{3. 开源软件更新}

项目组持续维护相关开源软件：\textbf{UnitDiskMapping.jl}（新增三角晶格编码模块）和\textbf{GadgetSearch.jl}（基于整数规划的小工具自动发现算法）。

\subsection*{与本申请项目的关系}

青年基金项目的核心进展（三角晶格编码的整数规划辅助搜索方法）直接构成本项目的方法论起点。本申请项目将在此基础上引入LLM驱动的规模化合成，从单条规约的半自动搜索扩展到百条规约的全自动合成，实现方法论的系统性提升。
